\section{Elementary forms of four auxiliary estimates}
\label{sec:elementary-inputs}

We give the versions of four classical estimates needed in the proof.
The Buchstab normalization follows directly from its delay equation;
the large-sieve and plane-curve bounds need only elementary integration
and elimination.  The progression upper bound is a consequence of
Lemma~\ref{lem:fundamental} and Mertens' product formula.  In particular,
the last argument retains the fundamental lemma as an input.

\subsection{The continuous Buchstab normalization}

The normalization used in Lemma~\ref{lem:future-integrals} is the
classical one; see \cite[Section~28.9, p.~236]{MontgomeryVaughanIII}.  The following
proof uses only the continuous delay equation, independently of the
positivity argument for the sieve functions.

\begin{lemma}\label{lem:elementary-buchstab}
Let $\omega:[1,\infty)\to\R$ be the continuous function determined by
\[
 u\omega(u)=1\quad(1\le u\le2),\qquad
 (u\omega(u))'=\omega(u-1)\quad(u>2).
\]
Then $\omega(u)\to e^{-\gamma}$ as $u\to\infty$, where
$\gamma=\lim_{n\to\infty}(H_n-\log n)$ and
$H_n=\sum_{j=1}^n1/j$.
\end{lemma}

\begin{proof}
Successive integration on unit intervals constructs a unique positive
continuous function, locally absolutely continuous and with bounded
one-sided derivatives on each such interval.  For $u>2$ its equation
gives, almost everywhere,
\begin{equation}\label{eq:elementary-buchstab-derivative}
 u\omega'(u)=\omega(u-1)-\omega(u)
             =-\int_{u-1}^{u}\omega'(v)\,dv.
\end{equation}
Put $M_n=\mathop{\rm ess\,sup}_{n<u<n+1}|\omega'(u)|$.
These quantities are finite by the construction, and $M_1\le1$.
For every integer $n\ge2$, \eqref{eq:elementary-buchstab-derivative}
implies
\[
 M_n\le\frac1n\max(M_{n-1},M_n).
\]
If $M_n>M_{n-1}$ this inequality forces $M_n=0$, a contradiction;
otherwise it gives $M_n\le M_{n-1}/n$.  Thus $M_n\le1/n!$.
In particular $\omega'$ is integrable on $[1,\infty)$, so $\omega$
is bounded and has a finite limit, say $\ell$.

It remains to determine $\ell$.  For $t>0$ define
\[
 W(t)=\int_1^\infty\omega(u)e^{-tu}\,du.
\]
Boundedness justifies differentiation and integration by parts.
The derivative of $u\omega(u)$ is zero on $(1,2)$; integrating it
on $[1,\infty)$, including the boundary value at $1$, gives
\[
 -e^{-t}-tW'(t)
   =\int_2^\infty\omega(u-1)e^{-tu}\,du
   =e^{-t}W(t).
\]
Since $W(t)\to0$ as $t\to\infty$, solving this differential equation
yields
\begin{equation}\label{eq:elementary-buchstab-transform}
 1+W(t)=\exp\!\left(\int_t^\infty\frac{e^{-s}}s\,ds\right).
\end{equation}
The constant in this formula can be evaluated without a Tauberian
argument.  Integration by parts shows that
\[
 \int_t^\infty\frac{e^{-s}}s\,ds+\log t
  =(1-e^{-t})\log t+\int_t^\infty e^{-s}\log s\,ds.
\]
Moreover,
\begin{equation}\label{eq:elementary-euler-integral}
 \int_0^\infty e^{-s}\log s\,ds=-\gamma.
\end{equation}
For an elementary verification, substitution $s=nv$ and integration
of the geometric sum give
\[
 \int_0^n(1-s/n)^n\log s\,ds
   =\frac{n}{n+1}\bigl(\log n-H_{n+1}\bigr).
\]
Indeed,
\[
 \int_0^1(1-v)^n\log v\,dv
  =-\frac1{n+1}\int_0^1\frac{1-(1-v)^{n+1}}v\,dv
  =-\frac{H_{n+1}}{n+1}.
\]
The integrands in the first identity, extended by zero for $s>n$,
are bounded in absolute value by $e^{-s}|\log s|$.  Dominated
convergence proves \eqref{eq:elementary-euler-integral}.
Consequently \eqref{eq:elementary-buchstab-transform} gives
$tW(t)\to e^{-\gamma}$ as $t\downarrow0$.
On the other hand, substituting $v=tu$ and using bounded convergence
gives
\[
 tW(t)=\int_t^\infty\omega(v/t)e^{-v}\,dv\longrightarrow\ell.
\]
This identifies $\ell=e^{-\gamma}$.
\end{proof}

This derivation supplies precisely the independently normalized limit
needed to rule out a jump at $2$ in the proof of
Lemma~\ref{lem:future-integrals}.  It makes no use of the auxiliary
Markov chain or of the future-integral identities proved there.

\subsection{The additive large sieve at the required order}

The sharp large sieve is due to Montgomery and Vaughan
\cite{MontgomeryVaughan1973}; see also \cite{VaughanLargeSieve}.
The following separated-arcs argument gives the order of
\eqref{eq:large-sieve-input} that is used in the inverse estimate.

\begin{lemma}\label{lem:elementary-large-sieve}
Let $x_1,\ldots,x_R\in\R/\Z$ have pairwise circular distance at least
$\delta$, where $0<\delta\le1$.  Let $J\ge1$ be an integer.  If complex
coefficients $b_j$ are supported on $J$ consecutive integers, then
\[
 \sum_{r=1}^{R}\left|\sum_j b_j e(jx_r)\right|^2
   \ll (J+\delta^{-1})\sum_j|b_j|^2
\]
with an absolute implied constant.  Hence
\eqref{eq:large-sieve-input} holds for every real $Q\ge1$.
\end{lemma}

\begin{proof}
Shifting the frequencies multiplies their sum by a complex number of
modulus one.  We may therefore write
$T(x)=\sum_{j=0}^{J-1}b_j e(jx)$ and $f(x)=|T(x)|^2$.
Let $I_r$ be the arc of length $\delta$ centred at $x_r$.
Their interiors are disjoint.  For $t\in I_r$, the fundamental theorem
of calculus along that arc gives
\[
 f(x_r)\le f(t)+\int_{I_r}|f'(x)|\,dx.
\]
Average over $t\in I_r$ and sum over $r$ to obtain
\[
 \sum_r f(x_r)
 \le\delta^{-1}\int_0^1|T(x)|^2\,dx
       +\int_0^1|f'(x)|\,dx.
\]
Since $|f'|\le2|T'||T|$, Cauchy--Schwarz and Parseval give
\[
 \int_0^1|f'(x)|\,dx
 \le2\|T'\|_2\|T\|_2
 \le4\pi(J-1)\sum_{j=0}^{J-1}|b_j|^2.
\]
This proves the asserted estimate.  Distinct reduced fractions modulo
one with denominators at most $Q$ have circular distance at least
$Q^{-2}$: subtracting two such fractions, or subtracting an integer
from their difference, leaves a nonzero integer numerator over the
product of their denominators.  Apply the estimate with
$\delta=Q^{-2}$.  When there is just one fraction, the same arc
argument applies directly.
\end{proof}

\subsection{A degree-uniform plane-curve intersection bound}

The appeals to B\'ezout's theorem in Lemma~\ref{lem:nonlinear-points}
require a bound uniform in both degrees and coefficients.  Here is the
needed affine form, proved by the elementary resultant argument.
Only distinct points are counted; multiplicities are unnecessary.

\begin{lemma}\label{lem:elementary-bezout}
Let $F,G\in\mathbb C[X,Y]$ be nonzero polynomials of respective total
degrees $m,n\ge1$, with no common nonconstant factor.  Then they have
at most $mn$ common zeros in $\mathbb C^2$.
\end{lemma}

\begin{proof}
First choose linear coordinates in which the coefficients of $Y^m$
in $F$ and $Y^n$ in $G$ are nonzero constants.  Such coordinates exist:
choose the direction of the $Y$-axis outside the zeros of the two
highest homogeneous parts.  Each part excludes only finitely many
directions in the projective line.  Scaling the polynomials makes
these leading coefficients one.

Put $K=\mathbb C(X)$.  The two polynomials are coprime in $K[Y]$.
Indeed, clearing denominators in a common factor and removing common
factors of its coefficients would give a common polynomial factor
in $\mathbb C[X,Y]$.  This is the elementary Gauss lemma: the product
of two polynomials whose coefficients have no common divisor again
has coefficients with no common divisor, as follows by reduction
modulo any irreducible divisor in $\mathbb C[X]$.

Consider the linear map over $K$
\begin{equation}\label{eq:elementary-sylvester-map}
 (A,B)\longmapsto AF+BG,
 \qquad \deg_Y A<n,\quad\deg_Y B<m,
\end{equation}
with target the polynomials of $Y$-degree less than $m+n$.
It is injective: from $AF=-BG$ and coprimality, $F$ divides $B$,
forcing $B=0$ and then $A=0$.  Its matrix in monomial bases is the
Sylvester matrix.  Its determinant $R(X)$ is a nonzero polynomial.
If $(x,y)$ is a common zero, every polynomial in the image of the
specialized map at $X=x$ vanishes at $Y=y$.  That image is a proper
subspace, so $R(x)=0$.  There are finitely many such $x$, and each
fiber contains finitely many points because $F(x,Y)$ is monic of
degree $m$.  Thus the common zero set is finite.

Now choose linear coordinates once more so that the $X$-coordinates
of those finitely many points are distinct and both leading
coefficients in $Y$ remain nonzero constants.  The first condition
excludes the finitely many directions joining pairs of common zeros;
the second excludes the directions on which either highest
homogeneous part vanishes.  A direction satisfying both exists.
Repeat the construction of $R$ in these coordinates.

We claim that $\deg R\le mn$.  Write
\[
 F=\sum_{k=0}^{m}f_k(X)Y^k,\qquad
 G=\sum_{k=0}^{n}g_k(X)Y^k,
 \qquad \deg f_k\le m-k,\quad\deg g_k\le n-k.
\]
View the rows of the Sylvester matrix as the coefficients of
$Y^iF$ for $0\le i<n$ and $Y^jG$ for $0\le j<m$; label its columns
by $Y^k$, $0\le k<m+n$.  An entry from row $Y^iF$ and column $Y^k$
has degree at most $m+i-k$; an entry from row $Y^jG$ has degree at
most $n+j-k$.  Every nonzero term of the determinant uses each row
and column once, so its degree is at most
\[
 \sum_{i=0}^{n-1}(m+i)+\sum_{j=0}^{m-1}(n+j)
       -\sum_{k=0}^{m+n-1}k=mn.
\]
Every common zero gives a distinct root of the nonzero polynomial
$R$.  This proves the bound.
\end{proof}

For clarity, the coprimality conditions in
Lemma~\ref{lem:nonlinear-points} can now be checked directly.
An irreducible factor $Q$ of degree $d\ge2$ that is not proportional
to a real polynomial is coprime to its coefficientwise conjugate.
Both partial derivatives of $Q$ are nonzero: independence of one
variable would force an irreducible polynomial over $\mathbb C$ to
be linear.  Their smaller degrees make them coprime to $Q$, giving
at most $d(d-1)$ intersections for each derivative; a nonzero constant
derivative has no common zeros with $Q$.  The same lemma
applies to the inflection polynomial
\eqref{eq:inverse-inflection} when it is nonzero and not divisible
by $Q$, since its degree is at most $3d-4$; a nonzero constant
inflection polynomial has no zeros.  If it is divisible by $Q$,
every nonspecial smooth real graph arc has second derivative zero
and lies on a line.  Restricting $Q$ to that line would then give a
univariate polynomial vanishing on an interval, hence identically;
the line would be a factor of $Q$.  This excludes such an arc and
is exactly the alternative used in that proof.

\subsection{An order bound for primes in a progression}

Brun--Titchmarsh gives a sharper constant than we need; see
\cite[Theorem~1.4.4]{GranvilleSoundararajan2014}.  We derive the
order bound in Lemma~\ref{lem:prime-inputs}(iii) from the fundamental
lemma, in the form recorded in \cite[Lemma~2.8]{Sofos2023}, and
Mertens' formula.  Working in progression coordinates keeps the
remainder independent of the modulus.

\begin{lemma}\label{lem:elementary-brun-titchmarsh}
There are absolute constants $C,J_0>0$ such that, for every integer
$q\ge1$, every integer $a$ with $(a,q)=1$, and every real $x$ and
$H>0$ with $J=H/q\ge J_0$,
\[
 \#\{x<p\le x+H:p\equiv a\pmod q\}
   \le C\frac{H}{\varphi(q)\log(H/q)}.
\]
\end{lemma}

\begin{proof}
Write the progression as $a+qj$, with $j$ in the real interval
\[
 (x-a)/q<j\le(x+H-a)/q
\]
of length $J$.  Fix an absolute $s\ge s_0$ sufficiently large that
$\theta=1/(2s)<1/2$, and put $v=J^\theta$ and $D=J^{1/2}=v^s$.
For every prime $p\le v$ not dividing $q$, divisibility of $a+qj$
by $p$ specifies exactly one forbidden class of $j$ modulo $p$.
For a squarefree product $l$ of these primes, the Chinese remainder
theorem therefore gives the intersection count
\[
 J/l+O(1),
\]
with an absolute error independent of $x,a,q,l$.
Apply Lemma~\ref{lem:fundamental} with $X=J$ and $g(p)=1/p$ on
this prime set.  If $N$ counts the surviving coordinates, it gives
\begin{equation}\label{eq:elementary-bt-sieve}
 N\ll J\prod_{\substack{p\le v\\p\nmid q}}(1-1/p)
       +D(1+\log D)^{C_2}
\end{equation}
for an absolute constant $C_2$.
To justify the remainder explicitly, let $C_1$ be the divisor-weight
exponent in \eqref{eq:fundamental-input} and choose an integer
$k\ge2^{C_1}$.  On squarefree $l$,
$\tau(l)^{C_1}\le\tau_k(l)$, where $\tau_k(l)$ counts ordered
factorizations of $l$ into $k$ positive integers.  Summing by the
first $k-1$ factors gives
\[
 \sum_{l\le D}\tau_k(l)
 \le D\left(\sum_{d\le D}\frac1d\right)^{k-1}
 \le D(1+\log D)^{k-1}.
\]
In particular, the remainder in \eqref{eq:elementary-bt-sieve}
has no dependence on $q$.

Mertens' formula gives, for large $J$,
\[
 \prod_{\substack{p\le v\\p\nmid q}}(1-1/p)
 =\prod_{p\le v}(1-1/p)
       \prod_{\substack{p\mid q\\p\le v}}(1-1/p)^{-1}
 \ll\frac1{\log v}\frac q{\varphi(q)}
 \ll\frac1{\log J}\frac q{\varphi(q)}.
\]
Every prime in the progression that exceeds $v$ is among the
survivors.  The primes at most $v$ contribute at most $v$ further
points, regardless of $q$.  Thus the desired prime count is at most
\[
 N+v\ll \frac{Jq}{\varphi(q)\log J}
          +J^{1/2}(1+\log J)^{C_2}+J^\theta.
\]
The last two terms are $O(J/\log J)$ for all sufficiently large
$J$, uniformly in $q$.  Since $q/\varphi(q)\ge1$, they are absorbed
by the first term.  Substituting $J=H/q$ completes the proof.
\end{proof}
